\section{Geodesic contours with several optical round trips}
\label{sec:optical-cycles}

The local cycle joins roots that both approach \(r=0\). To construct
a cycle with finite optical time, we use an outer root.
In addition to the \(m=d-2\) roots at
\(r=O(E^{-2/m})\), the turning-point equation has a pair of roots
asymptotic to \(\pm\ii E\). With \(r=E x\) the equation is
\(1+x^2-E^{-d}x^{-m}=0\), so the upper outer root is
\begin{equation}
 r_{\rm b}=\ii E[1+O(E^{-d})].
\end{equation}
For each fixed \(d\geq4\) and sufficiently large real \(E\), take the
path from \(r_0\) to the positive exterior through the upper horizon
detour of section~\ref{sec:radial}. Follow that leg to radius
\(R>2\max(E,1)\), rotate counterclockwise to \(\arg r_{\rm b}\), and
decrease the radius to \(r_{\rm b}\). On the exterior arc the
sheet is specified by \(Y\sim r\). Return from \(r_{\rm b}\)
to \(r_0\) on the other sheet. We denote the length and time
periods of this cycle by \(\Lambda_L\) and \(\Lambda_T\).
Every turn is regular at finite \(E\), since \(r-r_j\) is
quadratic in the local affine parameter at a simple root.

Insert \(n-1\) such cycles in the reference contour. The geodesic
then visits \(r_0\) \(n\) times and \(r_{\rm b}\) \(n-1\) times
before exiting on the reference outgoing sheet. For this sequence
at common energy the endpoint data are
\begin{align}
 L^{(n)}(E)&=L_0(E)+(n-1)\Lambda_L(E),&
 T^{(n)}(E)&=T_0(E)+(n-1)\Lambda_T(E).
 \label{eq:optical-insertions}
\end{align}
Only the two AdS endpoints require length counterterms. The repeated
turns take place on the complexified spacetime with the time lifted
along the contour.

\subsection{Optical time and finite-energy corrections}

The periods can be evaluated by subtracting twice the exterior
integral from \(r_{\rm b}\) to infinity from the reference result.
The exterior integral takes the ray out from \(r_{\rm b}\) and rotates
back to the positive axis. In this calculation the length subtraction
at the reference endpoint and at the exterior-integral endpoint
is the same. Writing \(r=r_{\rm b}y\), \(y\geq1\), gives
\begin{equation}
 \frac{Y^2}{r_{\rm b}^{\,2}}
 =y^2-1+r_{\rm b}^{-d}(1-y^{-m}) .
 \label{eq:optical-outer-factor}
\end{equation}
The ratio \((1-y^{-m})/(y^2-1)\) is bounded, including at the
turning point. Thus the doubled exterior length is
\(2\log(2/E)-\ii\pi+O(E^{-d})\), with the phase from transporting
the logarithm back to the positive axis. The leading integral
uses \(\int_1^Y (y^2-1)^{-1/2}\dd y=\operatorname{arcosh}(Y)\).
The leading exterior time comes from
\(\int_1^\infty [y^2\sqrt{y^2-1}]^{-1}\dd y=1\); multiplying by
\(2E/r_{\rm b}^{\,2}\) gives \(-2/E+O(E^{-d-1})\).

To resolve the first correction to the full cycle, compare
\(L_0(E)\) with the pure AdS half-line integral. Subtract their
integrands and write \(r=r_0 x\). Since \(r_0/E\sim E^{-q}\),
where \(q=d/(d-2)\), the coefficient of \(E^{-q}\) is
\begin{align}
 K_d
 &=2\left[\int_1^\infty
  \left(\frac{1}{\sqrt{1-x^{-m}}}-1\right)\dd x-1\right]
 =\frac{2}{m}B(-1/m,1/2).
 \label{eq:optical-kd}
\end{align}
Here the beta function is understood by continuation. The change
of variables \(v=x^{-m}\) subtracts the constant term in the
integrand at \(v=0\). The subtracted original integrand after
rescaling is bounded by a constant times \(x^{-m}\) for \(x\geq2\)
and by an integrable inverse square root near \(x=1\). The omitted
pure AdS interval \(0\leq r\leq r_0\) gives the \(-1\) in brackets.
Using \(L'_0=-E T'_0\) and the limiting \(T_*\) in
\eqref{eq:tstar} then gives
\begin{align}
 L_0(E)&=2\log(2/E)+K_d E^{-q}+o(E^{-q}),\notag\\
 T_0(E)&=T_*-\frac{2}{E}
         -\frac{q}{q+1}K_d E^{-q-1}+o(E^{-q-1}) .
 \label{eq:optical-ref-correction}
\end{align}
For \(d\geq4\), \(d>q\). Hence the leading exterior tail removes
the \(E^{-1}\) correction from each inserted cycle, and the
families in \eqref{eq:optical-insertions} have
\begin{align}
 L^{(n)}(E)&=2\log(2/E)+\ii\pi(n-1)+nK_d E^{-q}+o(E^{-q}),\notag\\
 T^{(n)}(E)&=nT_*-\frac2E
   -\frac{nq}{q+1}K_d E^{-q-1}+o(E^{-q-1}) .
 \label{eq:optical-period-asymptotics}
\end{align}
The \(1/E\) time correction comes from the two exterior endpoint
legs. In terms of \(\delta=nT_*-T^{(n)}\), elimination of \(E\)
gives
\begin{equation}
 L^{(n)}
   =2\log\delta+\ii\pi(n-1)
    +\frac{nK_d}{q+1}\left(\frac{\delta}{2}\right)^q
    +o(\delta^q).
 \label{eq:optical-fixed-time}
\end{equation}
Thus the classical worldline factor contains
\(e^{-\ii\pi M(n-1)}\delta^{-2M}\). Its contribution to a boundary
correlator also depends on the contour weight and fluctuation
determinant.

In AdS\(_5\), \(K_4=0\) and the cycle periods simplify exactly.
In \(y=r^2\) the doubled length integrand is
\(\dd y/\sqrt{y^2+E^2 y-1}\). The two cycle endpoints are the two
roots of this quadratic. The primitive
\(\operatorname{arcosh}[(2y+E^2)/\sqrt{E^4+4}]\) changes by
\(\ii\pi\) on the prescribed upper exterior path from the positive
root to the negative one. It follows that \(\Lambda_L=\ii\pi\).
The differential identity \(\Lambda'_L=-E\Lambda'_T\) and the
large-\(E\) limit then imply \(\Lambda_T=T_*\) for \(E>0\).

\begin{table}[t]
 \centering\small
 \begin{tabular}{ccccc}
 \hline
 \(d\) & \(n\) & \(E\) &
 \(E^q[L^{(n)}-2\log(2/E)-\ii\pi(n-1)]\) &
 \(E^{q+1}[T^{(n)}-nT_*+2/E]\) \\
 \hline
 5 & 2 & 8 & $-1.734206-0.002301i$ & $1.085715+0.001917i$ \\
5 & 2 & 32 & $-1.724900-0.000023i$ & $1.078094+0.000019i$ \\
5 & 2 & $\infty$ & $-1.724740$ & $1.077962$ \\
6 & 3 & 8 & $-3.600612+0.000000i$ & $2.161710+0.000000i$ \\
6 & 3 & 32 & $-3.594520+0.000000i$ & $2.156734+0.000000i$ \\
6 & 3 & $\infty$ & $-3.594421$ & $2.156652$ \\
8 & 3 & 8 & $-4.480939+0.000000i$ & $2.560517+0.000000i$ \\
8 & 3 & 32 & $-4.481005+0.000000i$ & $2.560574+0.000000i$ \\
8 & 3 & $\infty$ & $-4.481005$ & $2.560574$ \\
\hline

 \end{tabular}
 \caption{Finite-energy quadratures for geodesics with \(n\)
 visits to the inner turning point. The limit rows are the
 coefficients in \eqref{eq:optical-period-asymptotics}.}
 \label{tab:optical-cycles}
\end{table}

Table~\ref{tab:optical-cycles} compares two- and three-visit
contours with the analytic asymptotics. To check the contour
prescription, we compute the cycle directly over the two caps,
the explicit horizon detour and the exterior arc. The absolute
agreement with the exterior-tail subtraction is below
\(2\times10^{-14}\) in the examples shown. The scaled fixed-time
coefficient of \eqref{eq:optical-fixed-time} agrees with the
quadrature to absolute error below \(0.003\) at \(E=128\).
Replacing the upper by
the lower detour of the physical horizon changes the time by
\(\ii\beta\) per optical trip and leaves the length unchanged.
It therefore gives the limiting time \(nT_*^*\). In AdS\(_7\)
this is the time \(2n\mathcal Z\) of the primitive reflection
sectors in section~\ref{sec:einstein-poles}.
